\vspace{1.25in}

\question
Use the following phylogenetic tree to answer the following questions.
\begin{center}
\includegraphics[width=0.4\textwidth]{vertebrate_phylogeny}
\end{center}

\begin{parts}
\part[5]
Are the coelacanths more closely related to ray-finned fishes or to tetrapods. Tell how you know.

\begin{minipage}[t][0.9in]{0.9\textwidth}%

\begin{solution}
Coelacanths are more closely related to tetrapods because they share a more recent common ancestor with tetrapods than ray-finned fishes.
\end{solution}

\end{minipage}

\part[5]
Which group is a basal group (taxon)? What does this mean about it's evolution compared with that of the other groups? 

\begin{minipage}[t][0.9in]{0.9\textwidth}%
\begin{solution}
The jawless fishes are basal. That means these fishes evolved before the other groups of fishes and tetrapods.
\end{solution}
\end{minipage}

\part[3]
``Fishes'' includes all groups shown above except for tetrapods. Are fishes a monophyletic group? Tell how you know.

\begin{minipage}[t][0.9in]{0.9\textwidth}%
\begin{solution}
Fishes are not monophyletic. Fishes contain some but not all of the descendants of the common ancestor.
\end{solution}
\end{minipage}

\part[3]
What does the branch labeled ``1'' represent on this tree? What is the proper term for this branch?

%\begin{minipage}[t][0.6in]{\textwidth}%
\begin{solution}
The root, which represents the common ancestor of all groups shown on the phylogeny.
\end{solution}
%\end{minipage}

\end{parts}

\newpage

\question[14]
Population size tends to fluctuate around carrying capacity.  This ecological process can affect the evolutionary fitness of individuals in the population over time. The figure below shows the fluctuation of a frog population around \textit{K} over 20 years. Explain how competition for limited resources at points labeled A can increase the competitive ability of the frogs by natural selection as they repopulate from points labeled B over evolutionary time. (Hint: Remember how natural selection works.)%(Hint: think about the individuals that are most likely to survive during times when the population size is decreasing.)

%% THIS QUESTION CONTINUES TO NEED SOME WORK. TRY TO FIND A WAY TO EMPHASIZE THEY NEED TO EXPLAIN WHICH ARE SURVIVING AS POP SIZE DECREASES AND WHICH ARE REPOPULATING DURING INCREASE

\begin{center}
\includegraphics[width=0.5\textwidth]{fluctuating_population}
\end{center}

\begin{solution}
As the population size decreases, competition for limited resources is occurring. Those individuals that are better competitors for the resources will have the better chance of surviving and reproducing.  Those individuals that survive during the low points will, on average, have the better traits. As the population size increases, then the traits that allowed individuals to be better competitors will become more common in the population. By the time the population reaches point B, more individuals will have the beneficial trait.
\end{solution}